\begin{lemma}[Flat base change]
\label{lemma-flat-base-change-cohomology}
Consider a cartesian diagram of schemes
$$
\xymatrix{
X' \ar[d]_{f'} \ar[r]_{g'} & X \ar[d]^f \\
S' \ar[r]^g & S
}
$$
Let $\mathcal{F}$ be a quasi-coherent $\mathcal{O}_X$-module
with pullback $\mathcal{F}' = (g')^*\mathcal{F}$.
Assume that $g$ is flat and that $f$ is quasi-compact and quasi-separated.
For any $i \geq 0$
\begin{enumerate}
\item the base change map of
Cohomology, Lemma \ref{cohomology-lemma-base-change-map-flat-case}
is an isomorphism
$$
g^*R^if_*\mathcal{F} \longrightarrow R^if'_*\mathcal{F}',
$$
\item if $S = \Spec(A)$ and $S' = \Spec(B)$, then
$H^i(X, \mathcal{F}) \otimes_A B = H^i(X', \mathcal{F}')$.
\end{enumerate}
\end{lemma}

\begin{proof}
Using Cohomology, Lemma \ref{cohomology-lemma-base-change-map-flat-case} in (1)
is allowed since $g'$ is flat by
Morphisms, Lemma \ref{morphisms-lemma-base-change-flat}.
Having said this, part (1) follows from part (2). Namely,
part (1) is local on $S'$ and hence we may assume $S$
and $S'$ are affine. In other words, we have $S = \Spec(A)$
and $S' = \Spec(B)$ as in (2).
Then since $R^if_*\mathcal{F}$ is quasi-coherent
(Lemma \ref{lemma-quasi-coherence-higher-direct-images}),
it is the quasi-coherent $\mathcal{O}_S$-module associated to the
$A$-module $H^0(S, R^if_*\mathcal{F}) = H^i(X, \mathcal{F})$
(equality by
Lemma \ref{lemma-quasi-coherence-higher-direct-images-application}).
Similarly, $R^if'_*\mathcal{F}'$ is the quasi-coherent
$\mathcal{O}_{S'}$-module associated to the $B$-module
$H^i(X', \mathcal{F}')$. Since pullback by $g$ corresponds
to $- \otimes_A B$ on modules
(Schemes, Lemma \ref{schemes-lemma-widetilde-pullback})
we see that it suffices to prove (2).

\medskip\noindent
Let $A \to B$ be a flat ring homomorphism.
Let $X$ be a quasi-compact and quasi-separated scheme over $A$.
Let $\mathcal{F}$ be a quasi-coherent $\mathcal{O}_X$-module.
Set $X_B = X \times_{\Spec(A)} \Spec(B)$ and denote
$\mathcal{F}_B$ the pullback of $\mathcal{F}$.
We are trying to show that the map
$$
H^i(X, \mathcal{F}) \otimes_A B \longrightarrow H^i(X_B, \mathcal{F}_B)
$$
(given by the reference in the statement of the lemma)
is an isomorphism.

\medskip\noindent
In case $X$ is separated, choose an affine open covering
$\mathcal{U} : X = U_1 \cup \ldots \cup U_t$ and recall that
$$
\check{H}^p(\mathcal{U}, \mathcal{F}) = H^p(X, \mathcal{F}),
$$
see
Lemma \ref{lemma-cech-cohomology-quasi-coherent}.
If $\mathcal{U}_B : X_B = (U_1)_B \cup \ldots \cup (U_t)_B$ we obtain
by base change, then it is still the case that each $(U_i)_B$ is affine
and that $X_B$ is separated. Thus we obtain
$$
\check{H}^p(\mathcal{U}_B, \mathcal{F}_B) = H^p(X_B, \mathcal{F}_B).
$$
We have the following relation between the {\v C}ech complexes
$$
\check{\mathcal{C}}^\bullet(\mathcal{U}_B, \mathcal{F}_B) =
\check{\mathcal{C}}^\bullet(\mathcal{U}, \mathcal{F}) \otimes_A B
$$
as follows from
Lemma \ref{lemma-affine-base-change}.
Since $A \to B$ is flat, the same thing remains true on taking cohomology.

\medskip\noindent
In case $X$ is quasi-separated, choose an affine open covering
$\mathcal{U} : X = U_1 \cup \ldots \cup U_t$. We will use the
{\v C}ech-to-cohomology spectral sequence
Cohomology, Lemma \ref{cohomology-lemma-cech-spectral-sequence}.
The reader who wishes to avoid this spectral sequence
can use Mayer-Vietoris and induction on $t$ as in the proof of
Lemma \ref{lemma-quasi-coherence-higher-direct-images}.
The spectral sequence has $E_2$-page
$E_2^{p, q} = \check{H}^p(\mathcal{U}, \underline{H}^q(\mathcal{F}))$
and converges to $H^{p + q}(X, \mathcal{F})$.
Similarly, we have a spectral sequence with $E_2$-page
$E_2^{p, q} = \check{H}^p(\mathcal{U}_B, \underline{H}^q(\mathcal{F}_B))$
which converges to $H^{p + q}(X_B, \mathcal{F}_B)$.
Since the intersections $U_{i_0 \ldots i_p}$ are quasi-compact
and separated, the result of the second paragraph of the proof gives
$\check{H}^p(\mathcal{U}_B, \underline{H}^q(\mathcal{F}_B)) =
\check{H}^p(\mathcal{U}, \underline{H}^q(\mathcal{F})) \otimes_A B$.
Using that $A \to B$ is flat we conclude that
$H^i(X, \mathcal{F}) \otimes_A B \to H^i(X_B, \mathcal{F}_B)$
is an isomorphism for all $i$ and we win.
\end{proof}